\subsection{森田对应下保持的其他性质}

设 A 与 B 是森田等价的 k-代数，P 是一个可逆 \((A,B)_{k}\)-双模。

命题 11. —— 设

\[
\mathrm{V} ^ {\prime} \xrightarrow {f} \mathrm{V} \xrightarrow {g} \mathrm{V} ^ {\prime \prime}\tag{ε}
\]

是一个由 B-模与 B-线性映射组成的图，并设

\[
(\mathrm{P} \otimes \mathcal {E})
\]

\[
\mathrm{P} \otimes_ {\mathrm{B}} \mathrm{V} ^ {\prime} \xrightarrow {1 _ {\mathrm{P}} \otimes f} \mathrm{P} \otimes_ {\mathrm{B}} \mathrm{V} \xrightarrow {1 _ {\mathrm{P}} \otimes g} \mathrm{P} \otimes_ {\mathrm{B}} \mathrm{V} ^ {\prime \prime}
\]

是相应的 A-模的图。那么 \((\mathcal{E})\) 是正合列当且仅当 \((\mathrm{P} \otimes \mathcal{E})\) 是正合列。

设序列 \((\mathcal{E})\) 是正合的。由于右 B-模 P 是投射的，序列 \((P \otimes \mathcal{E})\) 是正合的（II, §3, No.~6, 第 251 页，命题 5 及 II, §3, No.~7, 第 257 页，推论 6）。

反之，设序列 \((\mathrm{P} \otimes \mathcal{E})\) 是正合的。设 Q 是一个 \((\mathrm{B}, \mathrm{A})_{k}\)-双模，它是 P 的逆，并设 \(\theta : Q \otimes_{A} P \to {}_{s} B_{d}\) 是一个同构。考虑交换图

\[
\begin{array}{c} \mathrm{Q} \otimes_ {\mathrm{A}} \mathrm{P} \otimes_ {\mathrm{B}} \mathrm{V} ^ {\prime} \xrightarrow {1 _ {\mathrm{Q}} \otimes 1 _ {\mathrm{P}} \otimes f} \mathrm{Q} \otimes_ {\mathrm{A}} \mathrm{P} \otimes_ {\mathrm{B}} \mathrm{V} \xrightarrow {1 _ {\mathrm{Q}} \otimes 1 _ {\mathrm{P}} \otimes g} \mathrm{Q} \otimes_ {\mathrm{A}} \mathrm{P} \otimes_ {\mathrm{B}} \mathrm{V} ^ {\prime \prime} \\ \theta \otimes 1 _ {\mathrm{V} ^ {\prime}} \Biggl \downarrow \\ \mathrm{V} ^ {\prime} \xrightarrow {f} \mathrm{V} \xrightarrow {g} \mathrm{V} ^ {\prime \prime}. \end{array}
\]

由于 Q 是投射 A-模且序列 \((\mathrm{P} \otimes \mathcal{E})\) 是正合的，此图的第一行是一个正合列。由于诸竖直箭头都是同构，第二行也是正合的。

推论. —— 设 \(f: V \rightarrow W\) 是一个 B-线性映射。那么 f 是单射（相应地，满射）当且仅当 \(1_{P} \otimes f\) 是。

命题 12. —— 设 V 是一个左 B-模。B-模 V 是投射的（相应地，是生成模、忠实的、*内射的、有限表现的*）当且仅当 A-模 P ⊗\_B V 亦然。

\begin{enumerate}
\def\labelenumi{\alph{enumi})}
\item
  设 V 是投射的。存在一个集合 I，使 V 同构于 \(\mathrm{B}_{s}^{(\mathrm{I})}\) 的一个直因子子模。于是 A-模 \(P \otimes_{B} V\) 同构于 \(\mathrm{P}^{(\mathrm{I})}\) 的一个直因子子模；由于 P 是投射 A-模，\(P \otimes_{B} V\) 亦然。
\item
  设 B-模 V 是生成模。设 M 是一个 A-模。存在一个 B-模 W，使 M 同构于 \(P \otimes_{B} W\)。由 VIII, 第 80 页的定理 1，存在一个集合 I 与一个满射 \(\varphi : V^{(I)} \to W\)。由该推论，映射 \(1_{P} \otimes \varphi\)（从 \(P \otimes (V^{(I)})\) 到 \(P \otimes_{A} W\)）是满射，这给出一个满射 \((P \otimes V)^{(I)} \to M\)。由 VIII, 第 80 页的定理 1，\(P \otimes V\) 是生成 A-模。
\item
  B-模 V 是忠实的当且仅当它的零化理想约化为 0。于是断言 c) 由 VIII, 第 105 页的例 3 得出。
\end{enumerate}

*d) 设 V 是内射的。由 VIII, 第 106 页的注，A-模 \(\mathrm{P} \otimes_{\mathrm{B}} \mathrm{V}\) 同构于 \(\mathrm{Hom}_{\mathrm{B}}(\mathrm{Q}, \mathrm{V})\)，其中 Q 是一个与 A 互逆的 \((\mathrm{B}, \mathrm{A})_k\)-双模。由于 A-模 Q 是投射的，因而是平坦的（X, §1, n° 3, 第 9 页，例 1），由 X, §1, n° 8, 第 18 页，命题 11，A-模 \(\mathrm{Hom}_{\mathrm{B}}(\mathrm{Q}, \mathrm{V})\) 是内射的。

\begin{enumerate}
\def\labelenumi{\alph{enumi})}
\setcounter{enumi}{4}
\item
  设 V 有一个有限表现 \(L_{1} \to L_{0} \to V \to 0\)（X, §1, n° 4, 第 10 页）。用 P 作张量积，我们导出 A-模的一个正合列 \(N_{1}^{\prime} \xrightarrow{u} N_{0}^{\prime} \to P \otimes_{B} V \to 0\)，其中 \(N_{1}^{\prime}\) 与 \(N_{0}^{\prime}\) 是投射的且有限生成的（命题 11 与 a)）。设 \(N_{0}^{\prime\prime}\) 是一个有限生成 A-模，使得模 \(N_{0} = N_{0}^{\prime} \oplus N_{0}^{\prime\prime}\) 是自由的且有限生成的，并设 \(u^{\prime}: N_{1}^{\prime} \oplus N_{0}^{\prime\prime} \to N_{0}\) 是同态 \((u, 1_{N_{0}^{\prime\prime}})\)；于是 \(P \otimes_{B} V\) 可以等同于 \(u^{\prime}\) 的余核。设 \(N_{1}\) 是一个有限生成的自由 A-模，\(p: N_{1} \to N_{1}^{\prime} \oplus N_{0}^{\prime\prime}\) 是一个满射同态；序列 \(N_{1} \xrightarrow{u^{\prime} \circ p} N_{0} \to P \otimes_{B} V \to 0\) 是 A-模 \(P \otimes_{B} V_*\) 的一个有限表现
\item
  设 A-模 \(P \otimes_{B} V\) 是投射的（相应地，是生成模、忠实的、*内射的、有限表现的*）。把上面的结果应用于（交换 A 与 B 以及 P 与 Q 的角色），可见 B-模 \(Q \otimes_{A} P \otimes_{B} V\) 也具有这一性质。于是对同构于它的 B-模 V 同样如此。
\end{enumerate}

推论. —— 环 A 是左阿廷的（相应地，左诺特的）当且仅当环 B 亦然。

由于 B 的左理想的有序集与 P 的 A-子模之集之间的同构，环 B 是左阿廷的（相应地，左诺特的）当且仅当 A-模 P 是阿廷的（相应地，诺特的）。而由 VIII, 第 101 页的定理 1，A-模 P 是生成模且有限生成；特别地，\(A_{s}\) 同构于 \(P^{n}\) 的一个直因子（对一个整数 \(n \geqslant 1\)）。因此，P 是阿廷的（相应地，诺特的）当且仅当 A 是左阿廷的（相应地，左诺特的）。

